% Part: second-order-logic
% Chapter: syntax-and-semantics

\documentclass[../../../include/open-logic-chapter]{subfiles}

\begin{document}

\olchapter{sol}{syn}{वाक्यविन्यास और अर्थविज्ञान}

\begin{editorial}
अब तक द्वितीय-कोटि तर्कशास्त्र का मूल वाक्यविन्यास और अर्थविज्ञान समाहित
किया गया है। एक अध्याय के रूप में यह अभी बहुत छोटा है। द्वितीय-कोटि चरों के
लिए प्रतिस्थापन पर चर्चा आवश्यक है, तभी द्वितीय-कोटि तर्कशास्त्र की
!!{derivation} प्रणालियों की बात की जा सकेगी; इसमें कुछ सूक्ष्म समस्याएँ भी हैं।
\end{editorial}

\olimport{introduction}

\olimport{terms-formulas}

\olimport{satisfaction}

\olimport{semantic-notions}

\olimport{expressive-power}

\olimport{inf-count}

\OLEndChapterHook

\end{document}
